% Part: incompleteness
% Chapter: arithmetization-syntax
% Section: proofs-in-ax

\documentclass[../../../include/open-logic-section]{subfiles}

\begin{document}

% SEGMENT OLP-0288-S01

\olfileid{inc}{art}{pax}
\olsection{அடிகோள் முறை \usetoken{P}{derivation}}

\begin{explain}
அடிகோள் முறை !!{derivation}s ஐ எண்கணிதமாக்க, !!{derivation}s ஐ
எண்களாகப் பிரதிநிதித்துவப்படுத்த வேண்டும். !!{derivation}s என்பவை
!!{formula}s இன் தொடர்கள் மட்டுமே என்பதால், ஒவ்வொரு !!{derivation} ஐயும்
அதிலுள்ள !!{formula}s இன் குறியீடுகளைக் கொண்ட தொடரின் குறியீடாகக்
குறியீடாக்குவதே தெளிவான அணுகுமுறை.
\end{explain}

\begin{defn}
$\delta$ என்பது $!A_1$, \dots,~$!A_n$ என்ற !!{formula}s கொண்ட ஓர்
அடிகோள் முறை !!{derivation} எனில், $\Gn{\delta}$ என்பது
\[
\tuple{\Gn{!A_1}, \dots, \Gn{!A_n}}.
\]
\end{defn}

\begin{ex}
  மிக எளிய பின்வரும் !!{derivation} ஐக் கருதுக:
  \begin{derivation}
  1. & $!B \lif (!B \lor !A)$ \\
  2. & $(!B \lif (!B \lor !A)) \lif (!A  \lif (!B \lif (!B \lor !A)))$\\
  3. & $!A  \lif (!B \lif (!B \lor !A))$
  \end{derivation}
  இந்த !!{derivation} இன் கோடல் எண்
  \begin{align*}
  \openTuple\,
    & \Gn{!B \lif (!B \lor !A)}, \\
    &\Gn{(!B \lif (!B \lor !A)) \lif (!A  \lif (!B \lif (!B \lor !A)))},\\
    & \Gn{!A  \lif (!B \lif (!B \lor !A))} \,\closeTuple.
  \end{align*}
  ஆக இருக்கும்.
\end{ex}

\begin{explain}
!!{derivation}s கான ஒரு பிரதிநிதித்துவத்தைத் தீர்மானித்தபின், அத்தகைய
!!{derivation}s ஐ முதன்மை மீள்வரையறை முறையில் கையாளலாம் என்றும், அவற்றின்
இன்றியமையாப் பண்புகளையும் தொடர்புகளையும் அவ்வாறே வெளிப்படுத்தலாம் என்றும்
காட்ட வேண்டும். சில செயல்கள் எளியவை: எடுத்துக்காட்டாக,
!!a{derivation} இன் கோடல் எண்~$d$ தரப்பட்டால்,
$(d)_{\len{d}-1}$ அதன் இறுதி-!!{formula} வின் கோடல் எண்ணைத் தருகிறது.
சில செயல்கள் மிகவும் கடினமானவை. அவற்றை எவ்வாறு செய்வது என்பதையாவது
வரைகோடாகக் காட்டுவோம். “$\delta$ என்பது~$\Gamma$ இலிருந்து~$!A$ இன்
!!a{derivation}” என்ற தொடர்பு, $\delta$ மற்றும்~$!A$ இன் கோடல் எண்களில்
முதன்மை மீள்வரையறையானது என்று காட்டுவதே இலக்கு.
\end{explain}

\begin{prop}
\ollabel{prop:followsby}
பின்வரும் தொடர்புகள் முதன்மை மீள்வரையறையானவை:
\begin{enumerate}
\item $!A$ ஓர் அடிகோள்.
\item $\delta$ வின் $i$-ஆம் வரி modus ponens மூலம் நியாயப்படுத்தப்படுகிறது.
\item $\delta$ வின் $i$-ஆம் வரி \QR{} மூலம் நியாயப்படுத்தப்படுகிறது.
\item $\delta$ ஒரு சரியான !!{derivation}.
\end{enumerate}
\end{prop}

\begin{proof}
!!{formula}s இன் கோடல் எண்களுக்கும் !!{derivation}s இன் கோடல் எண்களுக்கும்
இடையிலான இவற்றிற்கு ஒத்த தொடர்புகள் முதன்மை மீள்வரையறையானவை என்று காட்ட
வேண்டும்.
\begin{enumerate}
\item அடிகோள் வார்ப்புருக்களின் பட்டியல் ஒன்று நமக்குத் தரப்பட்டுள்ளது;
  அந்த வார்ப்புருக்களில் ஒன்றால் தரப்படும் வடிவில் $!A$ இருந்தால், அது ஓர்
  அடிகோள். வார்ப்புருக்களின் பட்டியல் முடிவுறுவது என்பதால், ஒவ்வோர் அடிகோள்
  வார்ப்புருவுக்கும் $!A$ அந்த வடிவில் உள்ளதா என்பதை முதன்மை மீள்வரையறை
  முறையில் சோதிக்கலாம் என்று காட்டினால் போதும். எடுத்துக்காட்டாக, பின்வரும்
  அடிகோள் வார்ப்புருவைக் கருதுக:
  \[
  !B \lif (!C \lif !B).
  \]
  ‘$($’ உடன் $!B$ ஐயும், அதனுடன் ‘$\lif$’ ஐயும், அதனுடன் ‘$($’ ஐயும்,
  அதனுடன் $!C$ ஐயும், அதனுடன் ‘$\lif$’ ஐயும், அதனுடன் $!B$ ஐயும்,
  இறுதியாக~‘$))$’ ஐயும் தொடரிணைத்தால் $!A$ கிடைக்குமாறு
  !!{formula}s $!B$, $!C$ இருந்தால், $!A$ இவ்வடிகோள் வார்ப்புருவின் ஒரு
  நேர்வு. $!A$ இன் கோடல் எண்~$n$ க்கு ஒத்த பண்பைச் சோதிக்கலாம்; ஏனெனில்
  தொடர்களின் தொடரிணைப்பு முதன்மை மீள்வரையறையானது; மேலும் $!B$, $!C$ உட்பட
  இவ்விரு வாய்பாடுகளின் கோடல் எண்கள்~$!A$ இன் கோடல் எண்ணைவிடச்
  சிறியவையாக இருக்க வேண்டும்; ஏனெனில் தொடர்பு பொருந்தும்போது $!B$, $!C$
  இரண்டும்~$!A$ இன் உட்-!!{formula}s. எனவே பின்வருமாறு வரையறுக்கலாம்:
  \begin{multline*}
  \fn{IsAx}_{!B \lif (!C \lif !B)}(n) \defiff \bexists{b< n}{
    \bexists{c < n}{(\fn{Sent}(b) \land \fn{Sent}(c) \land {}}}\\ n =
  \Gn{(} \concat b \concat \Gn{\lif} \concat \Gn{(} \concat c \concat
  \Gn{\lif} \concat b \concat \Gn{))}).
  \end{multline*}
  ஒவ்வோர் அடிகோள் வார்ப்புருவுக்கும் இத்தகைய வரையறை இருந்தால், அவற்றின்
  பிரிப்பிணைவு “$n$ ஓர் அடிகோளின் கோடல் எண்” என்ற $\fn{IsAx}(n)$
  பண்பை வரையறுக்கிறது.
\item $\delta$ வின் $i$-ஆம் வரி modus ponens மூலம் நியாயப்படுத்தப்படுவது,
  அப்போதும் அப்போது மட்டுமே, $j$, $k < i$ என்ற வரிகள் இருந்து,
  வரி~$j$ இலுள்ள !!{sentence} ஏதாவது வாய்பாடு~$!A$ ஆகவும், வரி~$k$
  இலுள்ள வாக்கியம் $!A \lif !B$ ஆகவும், வரி~$i$ இலுள்ள வாக்கியம்~$!B$
  ஆகவும் இருக்கும்.
  \begin{multline*}
    \fn{MP}(d, i) \defiff \bexists{j < i}{\bexists{k < i}{}}\\
    (d)_k = \Gn{(} \concat (d)_j
      \concat \Gn{\lif} \concat (d)_i \concat \Gn{)}
  \end{multline*}
  வரம்புள்ள அளவையிடல், தொடரிணைப்பு, மற்றும் $=$ ஆகியவை முதன்மை
  மீள்வரையறையானவை என்பதால், இது ஒரு முதன்மை மீள்வரையறைத் தொடர்பை
  வரையறுக்கிறது.
\item $\delta$ வின் ஒரு வரி $!B \lif \lforall[x][!A(x)]$ வடிவிலும்,
  அதற்கு முந்தைய ஒரு வரி ஏதாவது !!{constant}~$c$ க்கு
  $!B \lif !A(c)$ ஆகவும், $c$ ஆனது~$!B$ இல் இடம்பெறாமலும் இருந்தால், அந்த
  வரி \QR{} மூலம் நியாயப்படுத்தப்படுகிறது. இது பொருந்துவது, அப்போதும்
  அப்போது மட்டுமே,
  \begin{enumerate}
  \item ஒரு !!a{sentence}~$!B$ இருந்து,
  \item $x$ மட்டுமே கட்டிலா மாறியாகக் கொண்ட !!a{formula}~$!A(x)$ இருந்து,
  \item வரி~$i$ ஆனது $!B \lif \lforall[x][!A(x)]$ ஐக் கொண்டு,
  \item ஏதாவது மாறிலிக் குறிக்காக, ஏதாவது வரி $j < i$ ஆனது
    $!B \lif \Subst{!A}{c}{x}$ ஐக் கொண்டு,
  \item அந்த~$c$ ஆனது~$!B$ இல் இடம்பெறாமல் இருந்தால்.
  \end{enumerate}
  இவை அனைத்தையும் முதன்மை மீள்வரையறை முறையில் சோதிக்கலாம்; ஏனெனில்
  $!B$, $!A(x)$, $x$ ஆகியவற்றின் கோடல் எண்கள் வரி~$i$ இலுள்ள வாய்பாட்டின்
  கோடல் எண்ணைவிடச் சிறியவை; மேலும் $c$ இன் கோடல் எண் வரி~$j$ இலுள்ள
  வாய்பாட்டின் கோடல் எண்ணைவிடச் சிறியது:
  % SOURCE CORRECTION TA-INC-010: c, not a, is bounded by the formula on line j.
  % SOURCE CORRECTION TA-INC-011: bind the previously free preceding-line index j.
  \begin{multline*}
    \fn{QR}_1(d, i) \defiff \bexists{j<i}{\bexists{b < (d)_i}{\bexists{x <
        (d)_i}{\bexists{a < (d)_i}{\bexists{c < (d)_j}{(}}}}}
    \\ \fn{Var}(x) \land \fn{Const}(c) \land {} \\
    (d)_i = \Gn{(} \concat
    b \concat \Gn{\lif} \concat \Gn{\lforall} \concat x \concat a
    \concat \Gn{)} \land {}\\
    (d)_j = \Gn{(} \concat b \concat
    \Gn{\lif} \concat \fn{Subst}(a,c,x) \concat \Gn{)} \land {}\\ \fn{Sent}(b)
    \land \fn{Sent}(\fn{Subst}(a,c,x)) \land {} \bforall{k <
      \len{b}}{(b)_k \neq (c)_0})
  \end{multline*}
  இங்கு $c$, $x$ ஆகியவை முறையே மாறிலிக் குறியையும் மாறியையும்
  சொற்களாகக் கருதியவற்றின் கோடல் எண்கள் (அதாவது, அவற்றின் குறிக்
  குறியீடுகள் அல்ல) என்று எடுகொள்கிறோம். $\Subst{!A(x)}{c}{x}$ ஒரு
  !!a{sentence} ஆக உள்ளதா என்று சோதிப்பதன் மூலம் $x$ மட்டுமே
  $!A(x)$ இன் கட்டிலா மாறியா என்று சோதிக்கிறோம்; $!B$ இன் ஒவ்வொரு குறியும்
  $c$ இலிருந்து வேறுபட்டதாக இருக்க வேண்டும் என்று கோருவதன் மூலம்
  $c$ ஆனது~$!B$ இல் இடம்பெறவில்லை என்பதை உறுதிசெய்கிறோம்.

  \QR{} இன் மற்றொரு வடிவைப் பயிற்சியாக விட்டுவைக்கிறோம்.
\item $d$ ஆனது சரியான !!{derivation} ஒன்றின் கோடல் எண்ணாக இருப்பது,
  அப்போதும் அப்போது மட்டுமே, அதன் ஒவ்வொரு வரியும் ஓர் அடிகோளாகவோ,
  modus ponens அல்லது~\QR{} மூலம் நியாயப்படுத்தப்பட்டதாகவோ இருக்கும்.
  எனவே:
  \[
  \fn{Deriv}(d) \defiff \bforall{i < \len{d}}{(\fn{IsAx}((d)_i) \lor
    \fn{MP}(d,i) \lor \fn{QR}(d, i))}
  \]
\end{enumerate}
\end{proof}

\begin{prob}
\olref[inc][art][pax]{prop:followsby} இல் உள்ளவாறு பின்வரும்
தொடர்புகளை வரையறுக்க:
\begin{enumerate}
\item $\fn{IsAx}_{!A \lif (!B \lif (!A \land !B))}(n)$,
\item $\fn{IsAx}_{\lforall[x][!A(x)] \lif !A(t)}(n)$,
\item $\fn{QR}_{2}(d, i)$ (\QR{} இன் மற்றொரு வடிவுக்கு).
\end{enumerate}
\end{prob}

\begin{prop}
$\Gamma$ என்பது !!{sentence}s இன் முதன்மை மீள்வரையறைக் கணம் என்க.
அப்போது, “$x$ என்பது~$\Gamma$ இலிருந்து~$!A$ க்கான
!!a{derivation}~$\delta$ வின் குறியீடாகவும், $y$ என்பது~$!A$ இன்
கோடல் எண்ணாகவும் உள்ளது” என்பதை வெளிப்படுத்தும் $\Prf[\Gamma](x, y)$
என்ற தொடர்பு முதன்மை மீள்வரையறையானது.
\end{prop}

\begin{proof}
“$y \in \Gamma$” என்பது $R_\Gamma(y)$ என்ற முதன்மை மீள்வரையறைப்
பயனிலையால் தரப்படுகிறது என்க. $\Prf[\Gamma](x, y)$ என்ற தொடர்பு முதன்மை
மீள்வரையறையானது என்று காட்ட வேண்டும்; இதில் $\Prf[\Gamma](x, y)$ பொருந்துவது,
அப்போதும் அப்போது மட்டுமே, $y$ என்பது !!a{sentence}~$!A$ இன் கோடல்
எண்ணாகவும், $x$ என்பது~$\Gamma$ இலிருந்து~$!A$ க்கான
!!a{derivation} ஒன்றின் குறியீடாகவும் இருக்கும்.

முந்தைய கூற்றின்படி, $x$ என்பது சரியான !!{derivation}~$\delta$ ஒன்றின்
குறியீடாக இருந்தால், அப்போதும் அப்போது மட்டுமே பொருந்தும்
$\fn{Deriv}(x)$ என்ற பண்பு முதன்மை மீள்வரையறையானது. எனினும், அந்த
வரையறை, !!{derivation} இன் வரிகளை நியாயப்படுத்துவதற்கான கூடுதல் வழியாக
$\Gamma$ கணத்தை எடுத்துக்கொள்ளவில்லை. ஒரு வரி~\QR{} மூலம்
நியாயப்படுத்தப்படுகிறதா என்பதற்கான நமது முதன்மை மீள்வரையறைச் சோதனை,
மாறிலிக் குறி~$c$ ஆனது~$\Gamma$ வில் இடம்பெறக்கூடாது என்ற கோரிக்கையையும்
கருத்தில் கொள்ளவில்லை. $\Gamma$ வை நேரடியாகக் கருத்தில் கொள்ளுமாறு நமது
வரையறையைத் திருத்தலாம்; ஆனால் $\fn{Deriv}$ ஐயும் வருவித்தல் தேற்றத்தையும்
பயன்படுத்துவது எளிது. ஏதாவது முடிவுறு !!{sentence}s பட்டியல்
$!B_1$, \dots, $!B_n \in \Gamma$ இருந்து,
$\{!B_1, \dots, !B_n\} \Proves !A$ என்றால், அப்போதும் அப்போது மட்டுமே
$\Gamma \Proves !A$. மேலும் வருவித்தல் தேற்றத்தின்படி,
$\Proves (!B_1 \lif (!B_2 \lif \cdots (!B_n \lif !A)\cdots))$ எனில்
இது பொருந்தும். கோடல் எண்~$z$ கொண்ட !!a{sentence} ஒன்று இவ்வடிவில்
உள்ளதா என்பதை முதன்மை மீள்வரையறை முறையில் சோதிக்கலாம். ஆகவே,
கோடல் எண்~$y$ கொண்ட !!{sentence} இன்
\emph{$\Gamma$ விலிருந்து} !!a{derivation} ஒன்றின் கோடல் எண்ணாக~$x$ ஐக்
கருதுவதற்குப் பதிலாக, மேலுள்ள வடிவிலான ஒரு உள்ளமை நிபந்தனைவாக்கியத்தின்
$\emptyset$ இலிருந்தான !!a{derivation} இன் கோடல் எண்ணாக~$x$ ஐக்
கருதுகிறோம்.

முதலில், !!{sentence}s இன் ஒரு தொடர் நமக்கிருந்தால், அவை அனைத்தையும்
முன்னிலைகளாகவும் தரப்பட்ட !!{sentence} ஐப் பின்னிலையாகவும் கொண்ட
நிபந்தனைவாக்கியத்தை முதன்மை மீள்வரையறை முறையில் அமைக்கலாம்:
% SOURCE CORRECTION TA-INC-012: recurse through hCond, the three-argument helper.
% SOURCE CORRECTION TA-INC-013: use the stated primitive-recursive membership predicate.
\begin{align*}
  \fn{hCond}(s, y, 0) & = y \\
  \fn{hCond}(s, y, n+1) & = \Gn{(} \concat (s)_{n} \concat \Gn{\lif}
  \concat \fn{hCond}(s, y, n) \concat \Gn{)}\\
  \fn{Cond}(s, y) & = \fn{hCond}(s, y, \len{s})\\
  \intertext{எனவே $\Prf[\Gamma](x, y)$ ஐப் பின்வருமாறு வரையறுக்கலாம்:}
  \Prf[\Gamma](x, y) & \defiff \bexists{s < \fn{sequenceBound}(x,x)}{(} \\
  &\qquad (x)_{\len{x}-1} = \fn{Cond}(s,y) \land {} \\
  &\qquad\bforall{i<\len{s}}{R_\Gamma((s)_i)} \land {} \\
  &\qquad\fn{Deriv}(x)).
\end{align*}
ஒவ்வொரு $(s)_i$ உம் !!{derivation} இன் கடைசி வரியிலுள்ள
உட்-!!{formula} ஒன்றின் கோடல் எண், அதாவது $(x)_{\len{x}-1}$ ஐவிடச்
சிறியது என்று கருதுவதன் மூலம்~$s$ இன் வரம்பைப் பெறுகிறோம்.
முன்னிலைகள்~$!B \in \Gamma$ இன் எண்ணிக்கை, அதாவது~$s$ இன் நீளம்,
$x$ இன் கடைசி வரியின் நீளத்தைவிடச் சிறியது.
\end{proof}

\end{document}
